\chapter{Obsah přiloženého repozitáře}
\label{app:repository}

Zdrojové kódy, konfigurační soubory, experimentální skripty a~výsledky experimentu jsou součástí přiloženého repozitáře. Struktura repozitáře je popsána v~tabulce~\ref{tab:struktura-repozitare} v~kapitole~\ref{chap:metodika}; zde je uveden přehled hlavních komponent a~způsob reprodukce experimentu.

\section*{Hlavní komponenty}

\begin{description}
\item[\texttt{detection-mechanisms/}] Hlavní Python balíček obsahující implementaci všech jedenácti detektorů ze tří skupin (hluboké učení, klasické strojové učení, statistické), evaluační metriky, statistické testy a~vizualizace.
\item[\texttt{testing-environment/}] Kontejnerové testovací prostředí -- generátor syntetického provozu, přijímač a~agregátor toků, přehrávač srovnávacích dat.
\item[\texttt{experiments/}] Řízení experimentálního běhu (\texttt{run\_experiment.py}), výstupní metriky a~grafy v~\texttt{results/}.
\item[\texttt{tex/}] Zdrojové texty diplomové práce v~\LaTeX{u}.
\item[\texttt{Makefile}] Automatizace úloh \texttt{make} příkazy \texttt{install-all}, \texttt{prepare-data}, \texttt{evaluate-all}.
\end{description}

\section*{Reprodukce experimentu}

\begin{enumerate}
\item \texttt{make install-all} -- instalace závislostí.
\item \texttt{make prepare-data} -- stažení a~normalizace UNSW-NB15.
\item \texttt{make evaluate-all} -- spuštění experimentu pro všech 11~modelů včetně křížové validace.
\end{enumerate}

Iniciální hodnoty generátorů jsou fixovány na~42; výsledky jsou reprodukovatelné na shodné konfiguraci (Python~3.10, Linux).

\section*{Veřejné úložiště}

Zdrojové kódy, \LaTeX{} text práce i~experimentální data jsou veřejně dostupné na:
\begin{center}
\url{https://github.com/dashwoood/diploma}
\end{center}
